% Generated draft: requires mathematical and source review.
\documentclass[11pt]{article}
\usepackage[margin=25mm]{geometry}
\usepackage{fontspec}
\setmainfont{Latin Modern Roman}
\usepackage{amsmath,amssymb,mathtools,unicode-math}
\setmathfont{Latin Modern Math}
\usepackage{enumitem,needspace,xurl,hyperref}
\setlength{\parindent}{0pt}
\setlength{\parskip}{5pt}
\setlength{\emergencystretch}{3em}
\newcommand{\N}{\mathbb{N}}
\newcommand{\Z}{\mathbb{Z}}
\newcommand{\Q}{\mathbb{Q}}
\newcommand{\R}{\mathbb{R}}
\newcommand{\C}{\mathbb{C}}
\newcommand{\F}{\mathbb{F}}
\newcommand{\A}{\mathbb{A}}
\newcommand{\PP}{\mathbb{P}}
\newcommand{\E}{\mathbb{E}}
\newcommand{\eps}{\varepsilon}
\newcommand{\dd}{\,\mathrm{d}}
\begin{document}
\section*{UnsolvedMath ID: 30003885}
\textbf{Source problem code:} catalogue entry
\subsection{Definitions and mathematical statement}
The source does not specify the manifold category. The following gives a precise formulation for each fixed choice \(\mathcal C\in\{\text{smooth},\text{topological}\}\); the intended choice requires confirmation from the source. In the topological category, boundaries are required to be collared. These are conditional formulations, not an assertion that both categories are intended.

An integral homology \(3\)-sphere is a closed, connected, oriented \(3\)-manifold \(Y\) satisfying
\[
H_k(Y;\Z)\cong H_k(S^3;\Z)\qquad\text{for every }k\geq 0.
\]
Two such manifolds \(Y_0,Y_1\) are integrally homology cobordant in category \(\mathcal C\) if there is a compact, connected, oriented \(4\)-manifold \(W\) in that category with oriented boundary
\[
\partial W\cong (-Y_0)\sqcup Y_1,
\]
such that both boundary inclusions induce isomorphisms
\[
H_k(Y_i;\Z)\xrightarrow{\ \cong\ }H_k(W;\Z)
\qquad(i\in\{0,1\},\ k\geq0).
\]
Let \(\Theta^3_{\Z,\mathcal C}\) denote the abelian group of these cobordism classes, with addition induced by oriented connected sum, identity \([S^3]\), and inverse \(-[Y]=[-Y]\). This makes explicit the category suppressed in the source notation \(\Theta^3_{\Z}\).

For the category intended by the source, the two questions are:
\begin{enumerate}
\item Does there exist an injective group homomorphism
\[
\bigoplus_{j=1}^{\infty}\Z\longrightarrow\Theta^3_{\Z,\mathcal C}?
\]
Here the direct sum consists of integer sequences with finite support. Equivalently, do there exist classes \(x_1,x_2,\ldots\) such that, for every \(m\geq1\) and every \((a_1,\ldots,a_m)\in\Z^m\),
\[
\sum_{j=1}^{m}a_jx_j=0\quad\Longrightarrow\quad a_1=\cdots=a_m=0?
\]
\item Does there exist \(x\in\Theta^3_{\Z,\mathcal C}\setminus\{0\}\) and an integer \(n\geq2\) such that \(nx=0\)?
\end{enumerate}
The first question follows the explicit statements, which ask for a subgroup isomorphic to the displayed direct sum. The background assessment also uses the word ``summands.'' If a direct summand is intended, the stronger requirement is an isomorphism
\[
\Theta^3_{\Z,\mathcal C}\cong\left(\bigoplus_{j=1}^{\infty}\Z\right)\oplus B
\]
for some abelian group \(B\). The supplied record does not resolve this distinction.

The original statement describes these questions as open, while the supplied registry status is \texttt{partially\_solved}. Neither description is independently verified here.
\subsection{Short English statement}
For the intended category of integral homology cobordism, ask whether the group contains countably infinitely many independent elements of infinite order, and whether it contains a nonzero element of finite order. The category is unspecified, and the background leaves unclear whether the first question asks for a subgroup or a direct summand.
\subsection{Sources}
\begin{itemize}
\item Local UnsolvedMath registry, permanent ID \(30003885\), title ``Infinite Rank and Torsion in Integral Homology Cobordism,'' source problem code \texttt{catalogue entry}; supplied complete source record in the Oberwolfach Reports Open Problems collection. The registry lists proposed year \(2018\) and supplies no named proposer.
\item Supplied report citation: \emph{Topologie (2019)}, Oberwolfach Reports, pp.\ 1857--1910, DOI \texttt{10.4171/owr/2018/31}. \url{https://doi.org/10.4171/owr/2018/31}. The background describes these reports as records of mathematical workshops and their open-problem discussions. The citation's year \(2019\) is retained alongside the registry's proposed year \(2018\).
\item Supplied literature reference, labeled ``Manolescu, JAMS'': \url{https://doi.org/10.1090/jams829}. The supplied assessment attributes to Manolescu the disproof of the splitting condition underlying the high-dimensional triangulation conjecture and describes questions about torsion and infinite-rank summands in the homology-cobordism group as unresolved. This assessment is reproduced as supplied context, without independent verification or an inferred answer to either question.
\item Literature assessment date: \texttt{2026-08-21}. The supplied dated curation labels the record \texttt{partially\_solved} and lists ``Manolescu, JAMS'' and the Oberwolfach report as sources checked. This is supplied dated curation, not an independently verified literature assessment; the record itself requests independent verification of time-sensitive status and cited claims.
\item The supplied statement status is \texttt{corrected\_verified}; its accompanying verification note says that mathematical content was preserved while notation, typography, and omitted definitions were repaired for self-contained publication. This is a registry claim, not mathematical review of the present draft.
\end{itemize}
\end{document}
